164
\hfill
12 Surfaces and Differential Forms in $\mathbb{R}^n$
in it by any topological space homeomorphic to $\mathbb{R}^k$. Most often the standard param-
eter region for local charts is assumed to be an open cube $I^k$ or an open ball $B^k$
in $\mathbb{R}^k$. But this makes no substantial difference.

\noindent
To carry out certain analogies and in order to make a number of the following
constructions easier to visualize, we shall as a rule take a cube $I^k$ as the canonical
parameter domain for local charts on a surface. Thus a chart
\[
\varphi: I^k \to U \subset S
\]
(12.1)
gives a local parametric equation $x = \varphi(t)$ for the surface $S \subset \mathbb{R}^n$, and the $k$-
dimensional surface itself thus has the local structure of a deformed standard $k$-
dimensional interval $I^k \subset \mathbb{R}^n$.

\noindent
The parametric definition of a surface is especially important for computational
purposes, as will become clear below. Sometimes one can define the entire surface
by a single chart. Such a surface is usually called \textit{elementary}. For example, the graph
of a continuous function $f: I^k \to \mathbb{R}^{n-k}$ in $\mathbb{R}^n$ is an elementary surface. However,
elementary surfaces are more the exception than the rule. For example, our ordinary
two-dimensional terrestrial sphere cannot be defined by only one chart. An atlas of
the surface of the Earth must contain at least two charts (see Problem 3 at the end of
this section).

\noindent
In accordance with this analogy we adopt the following definition.

\noindent
\textbf{Definition 3} A set $\mathcal{A}(S) := \{\varphi_i: I^k \to U_i, i \in \mathbb{N}\}$ of local charts of a surface $S$
whose domains of action together cover the entire surface (that is, $S = \bigcup_i U_i$) is
called an \textit{atlas} of the surface $S$.

\noindent
The union of two atlases of the same surface is obviously also an atlas of the
surface.

\noindent
If no restrictions are imposed on the mappings (12.1), the local parametrizations
of the surface, except that they must be homeomorphisms, the surface may be situ-
ated very strangely in $\mathbb{R}^n$. For example, it can happen that a surface homeomorphic
to a two-dimensional sphere, that is, a topological sphere, is contained in $\mathbb{R}^3$, but
the region it bounds is not homeomorphic to a ball (the so-called Alexander horned
sphere).$^3$

\noindent
To eliminate such complications, which have nothing to do with the questions
considered in analysis, we defined a \textit{smooth} $k$-dimensional surface in $\mathbb{R}^n$ in Sect. 8.7
to be a set $S \subset \mathbb{R}^n$ such that for each $x_0 \in S$ there exists a neighborhood $U(x_0)$
in $\mathbb{R}^n$ and a diffeomorphism $\psi : U(x_0) \to \mathbb{R}^n = \{t \in \mathbb{R}^n \mid |t^i| < 1, i = 1, \dots, n\}$ un-
der which the set $U_S(x_0) := S \cap U(x_0)$ maps into the cube $I^k = \mathbb{R}^k \cap \{t \in \mathbb{R}^n \mid
t^{k+1} = \dots = t^n = 0\}$.

\noindent
It is clear that a surface that is smooth in this sense is a surface in the sense
of Definition 1, since the mappings $x = \psi^{-1}(t^1, \dots, t^k, 0, \dots, 0) = (\varphi^1(t), \dots, \varphi^k(t))$
\hfill 163
$^3$ An example of the surface described here was constructed by the American topologist
J.W. Alexander (1888–1977).